% Style based on the supplied Waseem_Generic_AI_ML_GenAI_AgenticAI.tex.
% Python replaces template tokens. Edit layout here, content in JSON/portfolio JS.
\documentclass[9pt,letterpaper]{extarticle}
\usepackage[margin=0.55in]{geometry}
\usepackage[T1]{fontenc}
\input{glyphtounicode}
\pdfgentounicode=1
\usepackage[utf8]{inputenc}
\usepackage{lmodern}
\usepackage{microtype}
\usepackage{pifont}
\usepackage{xcolor}
\usepackage{enumitem}
\usepackage{titlesec}
\usepackage[hidelinks]{hyperref}
\usepackage{tabularx,array,parskip,ragged2e,needspace}
\pagestyle{empty}
\setlength{\parindent}{0pt}
\setlength{\parskip}{1pt}
\definecolor{accent}{HTML}{155B8A}
\titleformat{\section}{\large\bfseries\color{accent}}{}{0pt}{}[\vspace{0.5pt}\color{accent}\titlerule]
\titlespacing*{\section}{0pt}{3pt}{2pt}
\setlist[itemize]{leftmargin=1.15em,itemsep=0.8pt,topsep=1pt,parsep=0pt,partopsep=0pt}
\newcommand{\role}[4]{%
  \Needspace{4\baselineskip}\textbf{#1}\hfill\textbf{#2}\\[-1pt]
  \textit{#3}\hfill\textit{#4}\par}
\newcommand{\entry}[2]{\Needspace{4\baselineskip}\textbf{#1}\hfill\textit{#2}\par}
\newcommand{\skillrow}[2]{\textbf{#1:} & #2\\}
\hypersetup{pdftitle={Waseem Raza, PhD - AI/ML for Wireless Networks},pdfauthor={Waseem Raza, PhD}}
\begin{document}
\begin{center}
{\LARGE\bfseries Waseem Raza, PhD}\\[2pt]
{\large\bfseries Machine Learning Engineer \textbar{} AI/ML for Wireless \& 5G/6G Systems}\\[1pt]
{\small Wireless AI \textbar{} 5G/6G \textbar{} RAN Intelligence \textbar{} Network Automation}\\[3pt]
{\footnotesize New Jersey, United States \quad \ding{37}\enspace \href{tel:+17328039703}{+1 732-803-9703} \quad \ding{41}\enspace \href{mailto:waseem.placements@gmail.com}{waseem.placements@gmail.com} \textbar{} \href{mailto:waseemraza844@gmail.com}{waseemraza844@gmail.com}\\[1pt]\href{https://www.linkedin.com/in/waseem-raza-phd-9a737858/}{LinkedIn} \textbar{} \href{https://github.com/WaseemRaza844}{GitHub} \textbar{} \href{https://scholar.google.com/citations?user=d58Drr0AAAAJ\&hl=en}{Google Scholar} \textbar{} \href{https://waseemraza844.github.io/portfolio/}{Portfolio}}
\end{center}
\vspace{-5pt}
\Needspace{3\baselineskip}
\section{Professional Summary}
AI/ML Engineer for Wireless Networks with a PhD in Electrical and Computer Engineering and 7+ years of combined industry and research experience in 5G/6G systems, network intelligence, propagation modeling, mobility, RAN optimization, anomaly detection, and resilient automation. Apply deep learning, transfer learning, domain adaptation, generative modeling, reinforcement learning, and statistical methods to network-assurance, NR-NTN, CSI compression, digital-twin, outage-management, and self-healing problems. Experienced with Python, PyTorch, TensorFlow, scikit-learn, SQL, large-scale network measurements, 3GPP-compliant simulation, and production-oriented ML validation.
\Needspace{3\baselineskip}
\section{Technical Skills}
\Needspace{2\baselineskip}\noindent{\small\textbf{AI/ML for Wireless:} Transfer learning, domain adaptation, adversarial learning, generative modeling, reinforcement learning, anomaly detection, classification, and robust model evaluation}\par
\Needspace{2\baselineskip}\noindent{\small\textbf{Wireless Systems:} LTE/4G, 5G-NR, NR-NTN, mobility, RAN, propagation, cell selection, user timing, network measurements, and service assurance}\par
\Needspace{2\baselineskip}\noindent{\small\textbf{Network Intelligence:} Root-cause analysis, outage detection and compensation, self-healing, optimization, resilience, and intelligent automation}\par
\Needspace{2\baselineskip}\noindent{\small\textbf{Simulation \& Digital Twins:} 3GPP-compliant simulation, propagation modeling, synthetic data generation, digital twins, large-scale experimentation, and robustness testing}\par
\Needspace{2\baselineskip}\noindent{\small\textbf{Programming \& Data:} Python, SQL, PyTorch, TensorFlow, scikit-learn, XGBoost, Pandas, NumPy, visualization, and statistical analysis}\par
\Needspace{2\baselineskip}\noindent{\small\textbf{Engineering \& MLOps:} Git, Docker, CI/CD, reproducible experimentation, APIs, cloud analytics, model validation, and production-oriented workflows}\par
\Needspace{4\baselineskip}\textbf{Leadership and Professional Strengths}\par
\begin{itemize}
\item \textbf{Ownership and accountability:} Lead AI/ML initiatives through requirements analysis, data preparation, model development, validation, stakeholder review, and deployment planning.
\item \textbf{Technical leadership:} Define technical approaches, evaluate alternatives, coordinate experimentation, and communicate findings to guide projects from ambiguous requirements to measurable outcomes.
\item \textbf{Cross-functional collaboration:} Work with software engineers, network specialists, researchers, geologists, and public-safety stakeholders to translate domain requirements into practical data and AI solutions.
\item \textbf{Continuous learning:} Apply ongoing study in generative AI, agentic AI, RAG, LLM evaluation, and cloud-native platforms through hands-on projects and professional coursework.
\item \textbf{Reliability and commitment:} Maintain reproducible workflows, careful validation, and attention to detail throughout experimentation and delivery.
\item \textbf{Adaptability and problem solving:} Combine critical thinking, systematic experimentation, and data-driven analysis to address unfamiliar technologies and complex operational problems.
\end{itemize}
\Needspace{3\baselineskip}
\section{Professional Experience}
\role{Machine Learning Engineer}{02/2025 - Present}{AT\&T, contractor via Innovatix Inc.}{Bedminster, NJ}
\begin{itemize}\interlinepenalty=10000
\item Lead AI/ML development for E911 and network-assurance diagnostics, applying NLP, domain adaptation, and root-cause classification to operational telecom workflows.
\item Develop multimodal diagnostic models using work-order text, location information, network context, and operational signals for VoNR and 5G service assurance.
\item Design cross-domain validation and distribution-shift evaluation workflows to improve reliability of ML models across heterogeneous network environments.
\item Collaborate with network engineers, software teams, and public-safety stakeholders on AI-enabled automation for critical telecommunications operations.
\end{itemize}
\role{Advanced-Degrees Intern}{06/2024 - 08/2024}{AT\&T Labs}{Austin, TX}
\begin{itemize}\interlinepenalty=10000
\item Analyzed 500K+ LTE/5G-NR diagnostic logs and applied ML methods to detect anomalies, reducing diagnostic time by 20\%.
\item Investigated 5G NR-NTN cell-selection behavior and improved decision algorithms by 15\% through data-driven analysis and validation.
\item Conducted structured analysis of large operational datasets and communicated technical findings to engineering teams supporting next-generation wireless systems.
\end{itemize}
\role{Advanced Wireless Technology Development Intern}{09/2022 - 01/2023}{Futurewei Technologies}{Schaumburg, IL}
\begin{itemize}\interlinepenalty=10000
\item Improved autoencoder-based channel-state-information compression by 25\% using transfer learning and systematic experimentation on large wireless datasets.
\item Conducted 50+ experiments to validate model performance, robustness, and generalization under changing data conditions.
\item Collaborated with research and engineering stakeholders to translate modeling results into recommendations for next-generation communication systems.
\end{itemize}
\role{Advanced-Degrees Intern}{06/2022 - 08/2022}{AT\&T Labs}{Bedminster, NJ}
\begin{itemize}\interlinepenalty=10000
\item Designed ML and statistical methods for analyzing users' timing in 5G standalone networks, improving analytical precision by 20\%.
\item Analyzed large-scale user trace data with SQL and visualization tools and communicated findings to senior technical leadership.
\item Converted raw operational traces into interpretable performance indicators to support diagnosis and network optimization.
\end{itemize}
\role{Graduate Research Assistant}{09/2019 - 07/2025}{AI4Networks, University of Oklahoma}{Tulsa, OK}
\begin{itemize}\interlinepenalty=10000
\item Designed reusable Python simulation and data-generation platforms for large-scale experimentation, including a 3GPP-compliant network simulator and an AI-enabled 5G testbed.
\item Built transfer-learning, adversarial-learning, GAN, autoencoder, and reinforcement-learning pipelines for sparse-data and cross-domain modeling.
\item Developed robustness evaluation frameworks for models operating under distribution shifts, missing data, and limited labeled samples; led 100+ experiments from problem formulation through validation.
\item Developed scalable self-healing and outage-management frameworks that improved modeled network resilience by 30\%.
\item Published peer-reviewed research, presented technical results internationally, and mentored junior researchers in ML implementation, experimental design, and technical writing.
\end{itemize}
\Needspace{3\baselineskip}
\section{Education}
\entry{PhD, Electrical and Computer Engineering}{09/2019 - 07/2025}
University of Oklahoma, Norman, Oklahoma, USA\par
\entry{MSc, Telecommunication Engineering}{08/2014 - 08/2016}
University of Engineering and Technology, Taxila, Pakistan\par
\entry{BSc, Telecommunication Engineering}{11/2010 - 06/2014}
University of Engineering and Technology, Taxila, Pakistan\par
\Needspace{3\baselineskip}
\section{Selected Projects}
\entry{SyntheticNET: 3GPP-Compliant Network Simulation}{2020-2024}
\textit{Research Project}\par
\begin{itemize}\interlinepenalty=10000
\item Python-based simulation environment for realistic mobility, propagation, and data-driven experimentation in cellular networks.
\end{itemize}
{\small\textit{Tools \& methods:} Python, 3GPP, Simulation, Wireless Networks\par}
\entry{TurboRAN Experimental Platform}{2020-2024}
\textit{Research Infrastructure}\par
\begin{itemize}\interlinepenalty=10000
\item Multi-band, multi-tier experimental platform supporting evaluation of emerging AI-enabled radio-access-network solutions.
\end{itemize}
{\small\textit{Tools \& methods:} 5G Testbeds, RAN, Experiment Design, Network Measurement\par}
\Needspace{3\baselineskip}
\section{Certifications}
\Needspace{3\baselineskip}\noindent\textbf{1. AI for Telecommunications Specialization} {\small --- \textcolor{blue}{\underline{\href{https://coursera.org/verify/specialization/D8Y6MVC0XM51}{AI CERTs / Coursera}}} \textbar{} Completed (3/3) courses \textbar{} 06/2026 \textbar{} \textcolor{blue}{\underline{\href{https://waseemraza844.github.io/portfolio/certificates/wireless/ai-for-telecom.pdf}{Certificate}}}}\par
\textbf{Focus:} Three-course specialization connecting AI foundations with telecommunications network optimization, security, and customer experience.\par
{\small\textbf{Tools \& Skills:} AI for Telecommunications, Network Optimization, Cybersecurity, Customer Experience, Internet of Things, Python, TensorFlow\par}
\vspace{2pt}
\Needspace{3\baselineskip}\noindent\textbf{2. 4G Network Fundamentals} {\small --- \textcolor{blue}{\underline{\href{https://coursera.org/verify/DLW1YN474VZJ}{Institut Mines-Télécom / Coursera}}} \textbar{} Completed (1/1) courses \textbar{} 08/2026 \textbar{} \textcolor{blue}{\underline{\href{https://waseemraza844.github.io/portfolio/certificates/wireless/4g-network-fundamentals.pdf}{Certificate}}}}\par
\textbf{Focus:} Foundational study of fourth-generation cellular networks.\par
{\small\textbf{Tools \& Skills:} 4G Networks, Cellular Networks, Telecommunications\par}
\vspace{2pt}
\Needspace{3\baselineskip}\noindent\textbf{3. Business Considerations for 5G with Edge, IoT, and AI} {\small --- \textcolor{blue}{\underline{\href{https://coursera.org/verify/SAUEVQXV5W71}{EDUCBA / Coursera}}} \textbar{} Completed (1/1) courses \textbar{} 02/2026 \textbar{} \textcolor{blue}{\underline{\href{https://waseemraza844.github.io/portfolio/certificates/wireless/bus-cons-5g.pdf}{Certificate}}}}\par
\textbf{Focus:} Explores business considerations for combining 5G connectivity with edge computing, IoT, and AI.\par
{\small\textbf{Tools \& Skills:} 5G, Edge Computing, Internet of Things, Artificial Intelligence, Business Analysis\par}
\vspace{2pt}
\Needspace{3\baselineskip}\noindent\textbf{4. Machine Learning} {\small --- Stanford University / Coursera \textbar{} Completed}\par
\textbf{Focus:} Core supervised and unsupervised learning, model evaluation, optimization, and practical ML system development.\par
{\small\textbf{Tools \& Skills:} Supervised Learning, Unsupervised Learning, Model Evaluation, Optimization\par}
\vspace{2pt}
\Needspace{3\baselineskip}
\section{Selected Publications}
\entry{On Multi-Parameter Optimization and Proactive Reliability in Cellular Networks}{2025}
\textit{Journal article}\par
\entry{An AI-Driven Framework for Enhancing Resilience in Propagation Models to Enable Digital Twin}{2024}
\textit{IEEE PIMRC}\par
\entry{Data-Driven Intelligent Outage Management for High Shadowing Environments in 5G\&B Networks}{2024}
\textit{IEEE SmartNets}\par
\entry{\href{https://arxiv.org/abs/2203.15899}{Toward a Hybrid RF/Optical Lunar Communication System (LunarComm)}}{2022}
\textit{IEEE Network}\par
\entry{Towards Positioning Error Impact Characterization and Minimization in User-Centric RAN}{2022}
\textit{IEEE WCNC}\par
\end{document}
